\resourceheading{9}{Full-Week Personalized Family Recap}{Family communication, student authorship, and learning continuity\par}{Students, families, teachers, and support teams}{Module 3 Daily Communication Logs critique and completed EDU486 day-by-day recap books}

\vspace{-0.5cm}
\toolsection{The Existing Camp Artifact}

This resource grew directly from how I wanted families to see the week: as a record of what they did, noticed, questioned, contributed, and helped shape. For each of our two assigned youth, we created a personalized full-week recap book. One completed book contains five landscape pages for Days 1--5, and each attended day has its own page built around an action image, a transcript-grounded contribution, a concise science story, and five activity memories connecting what the learner did to why it mattered \citep{campEvidence2026}. The science moves (observing, modeling, testing, asking, and acting) help make participation visible with dignity. Each page also includes an evidence boundary so that what we observed remains separate from what we interpreted or could not confirm. This gives learners and families something concrete they can point to, retell, question, or correct together. Personalized email drafts were also prepared to introduce the books and invite families to review them with the learner.

\toolsection{Per Day and Per Activity: The Completed Page}
\begin{center}

\restrictedartifact
  {assets/artifact-previews/weekly-recap-day3-example.png}
  {0.85\textwidth}

\parbox{1.0\textwidth}{\centering
\textbf{\color{DeepNavy}Personalized Day 3 Science Story.}

The completed page centers the learner as an investigator, preserves a grounded contribution, and reconstructs the day through activity memories. Each card connects action to thinking while the evidence boundary states what the camp record establishes. It functions as a learning record and a memory support: the learner can point to an activity, retell in their own way, and revisit how their thinking developed.
}

\end{center}

\newpage
\small

\toolsection{From Rating Log to Family Learning Story}

The EDE448 Daily Communication Logs packet offers useful structures for home--school communication, but several examples foreground good/fair/poor days, warnings, sitting, compliance, or adult initials \citep{dailylogs2009}. The completed camp recap books demonstrate a different communication architecture: they organize the record around what the learner did, communicated, questioned, contributed, and was still working to understand.

\renewcommand{\arraystretch}{1.00}
\setlength{\tabcolsep}{3pt}

\begin{tabularx}{\textwidth}{
>{\bfseries\RaggedRight\arraybackslash}p{1.19in}
>{\RaggedRight\arraybackslash}p{2.70in}
>{\RaggedRight\arraybackslash}X}

\rowcolor{MidTeal}
\color{white}Adult-Centered Pattern &
\color{white}Recap-Book Shift &
\color{white}Support Value \\

Good/fair/poor day &
Observable activity, contribution, learner words or action, and an evidence limit &
Context replaces moral rating \\

\rowcolor{SoftGray}
Warnings, sitting, compliance &
Choices, questions, refusals, sensory boundaries, roles, movement, and re-entry &
The learner appears as a communicator and participant, not a management problem \\

One summary box &
One page per attended day and five activity cards within each page &
Differences across activities, roles, and access needs remain visible \\

\rowcolor{SoftGray}
Adult initials as authorship &
Student words, artifacts, documented actions, learner review, and family invitation &
Authorship and correction are distributed rather than owned by one adult \\

Finished-sounding claim &
``Still to test,'' ``not confirmed,'' or another precise evidence boundary &
Uncertainty remains visible, honest, and useful \\

\rowcolor{SoftGray}
Required family return &
Optional response through a chosen mode, time, or not at all &
Family participation is invited without treating silence as lack of care \\

\end{tabularx}

\toolsection{Ready-to-Use Record / Student-Centered Entry (Reconstructed Example)}

\renewcommand{\arraystretch}{1.12}
\setlength{\tabcolsep}{4pt}
\begin{tabularx}{\textwidth}{>{\bfseries\RaggedRight\arraybackslash}p{1.55in} >{\RaggedRight\arraybackslash}X}
\rowcolor{MidTeal}
\color{white}Entry Field & \color{white}De-identified Camp-to-Home Example \\
Day and attendance & Day 3 attended (Invisible Invaders weekly recap sequence; five activity cards reviewed) \\
\rowcolor{SoftGray}
Observable participation & Joined sampling and evidence-talk tasks; pointed to sample locations, verbally compared findings with a partner, and used ``pass + return'' before rejoining group reporting \\
Access supports used & Visual activity sequence, point/tell/show/partner/pass + return routes, processing time, movement re-entry option, and choice to observe-only for high-smell materials \\
\rowcolor{SoftGray}
Concise camp-to-home update & ``Today your child helped compare sample evidence, named one hands-on part they liked, and returned to the group after a short break. The team marked what is confirmed vs. still to test in the recap.'' \\
Optional family-return area (blank by design) & \textbf{Optional reply only if helpful:} Home connection noticed: \underline{\hspace{1.85in}}\\[0.05in]
& Preferred support for next session: \underline{\hspace{1.5in}}\\[0.05in]
& Best return mode (text/email/paper/audio/none): \underline{\hspace{1.25in}} \\
\end{tabularx}

\vspace{0.04in}
\noindent\footnotesize\textit{Evidence boundary: Day attendance structure, activity-sequence reporting, access routes, and prepared family-invite workflow are documented in the camp archive; the blank family-return prompts above are a proposed reusable format and are not a recovered parent reply} \citep{campEvidence2026}.
\normalsize

\toolsection{Prepare, Review, and Deliver}

\begin{enumerate}
  \item \textbf{Verify identity and access.} Confirm attendance, name, pronouns, languages, intended readers, and permissions.
  \item \textbf{Assemble the evidence.} Gather student memos, artifacts, recap materials, images, and the activity sequence.
  \item \textbf{Separate evidence from interpretation.} Mark exact words, documented actions, team evidence, interpretations, and unknowns before designing the page.
  \item \textbf{Return control to the learner.} When possible, let the learner select, edit, translate, restrict, replace, or withhold parts of the outgoing story.
  \item \textbf{Review the complete record.} Check every attended day, all activity cards, attribution, accessibility, privacy, and evidence limits before export.
  \item \textbf{Deliver without overclaiming.} Use an approved communication channel, document receipt when available, and invite a low-burden response. Without evidence of delivery, say \textit{prepared for family delivery}, not \textit{sent}.
\end{enumerate}

\toolsection{Puzzle Plan and EQUITAS in Practice}

The Puzzle Plan carries the student memo, daily souvenir, team recap, activity evidence, and week-long storyline forward without collapsing them into one adult judgment. Each source contributes a different piece of the learner's experience. EQUITAS asks who controls the family narrative, whose activities or communication routes are missing, whether interpretation is clearly labeled, whether uncertainty remains visible, and whether the family has an accessible way to return knowledge to the team. Family knowledge is part of the evidence system \citep{zacarian2015community}.

\reflectionbox{\textbf{Artifact and delivery proof:} The private production bundle contains two final personalized recap PDFs, plus review renders, source records, image manifests, and two personalized email drafts \citep{campEvidence2026}.}

\adaptationbox{Family communication can become surveillance or burden. Collect only information that supports learning, communication, access, or continuity; follow privacy policy; use approved channels; and honor preferences for translation, audio, paper, event-triggered updates, weekly summaries, no photographs, no specific message, or no formal log at all.}

% \sourcebox{\scriptsize Daily Communication Logs \citep{dailylogs2009}; family recap archive \citep{campEvidence2026}; classroom-community partnership \citep{zacarian2015community}.}

\normalsize